\chapter{Anatomy of a manifest}
\label{ch:manifest}

The manifest is the single hand-written file that drives rosetta. This chapter
covers the rules that apply to \emph{all} of it --- how paths resolve, how to
write comments, how to factor out repetition --- before the following chapters
go field by field.

It answers three questions for the framework:

\begin{enumerate}
  \item \textbf{Where} are your headers, and where is rosetta's \code{include/}?
  \item \textbf{What} classes and free functions should be bound?
  \item \textbf{Which} language backends should be emitted?
\end{enumerate}

Everything else --- fields, methods, constructors, enums, inheritance --- is
discovered by reflection from the headers themselves.

\section{The minimum}

Three keys are required: \field{user\_include}, \field{rosetta\_include} and
\field{targets}, plus \field{classes} unless a preset supplies it.

\begin{lstlisting}[style=json]
{
  "user_include": "./geom",
  "rosetta_include": "../../include",
  "targets": ["python"],
  "classes": [ { "header": "Point.h" } ]
}
\end{lstlisting}

That is a complete, working manifest.

\section{Paths resolve from the manifest}

Every relative path in the file is resolved against \textbf{the directory
holding the manifest}, not the current working directory and not the output
tree. Absolute paths are taken as written.

\begin{warn}
Move the manifest and you must re-run \code{rosetta\_gen}. Nothing detects a
stale resolution for you.
\end{warn}

There is one deliberate exception, and it is the one people trip over: a class
entry's \field{header} is resolved against \field{user\_include}, not against
the manifest. That is because the header is emitted into an
\code{\#include "\ldots"} line in generated code --- it must be spelled relative
to the \emph{include path}, exactly as the compiler will look for it.
\field{user\_sources}, which name files to compile rather than to include, do
resolve from the manifest.

\section{Comments}

JSON has no comment syntax, so the manifest uses two conventions, and both work
everywhere:

\begin{lstlisting}[style=json]
{
  "//": "Bindings for the geom library.",
  "//paths": "All relative paths resolve from THIS file's directory.",

  // Line comments are tolerated by the parser too.
  /* As are block comments. */
  "user_include": "./geom"
}
\end{lstlisting}

Any key beginning with \code{//} is ignored --- \code{"//1"}, \code{"//note"},
\code{"//cpp26"}. Keys must still be unique, which is why the numbered and
named forms exist. The \code{//} and \code{/* */} line forms are accepted by
the parser as an extension; the key form is plain JSON and survives any tool
that reads the file.

\section{Variables}
\label{sec:variables}

A path that appears in several fields would otherwise be written out in each of
them. The \field{cpp26\_*} block is the case that motivated this --- one
toolchain root, spelled four times with a different suffix:

\begin{lstlisting}[style=json]
"cpp26_root": "$ENV{HOME}/devs/c++/clang-p2996/build",
"cpp26_cxx":  "$ENV{HOME}/devs/c++/clang-p2996/build/bin/clang++",
"cpp26_cc":   "$ENV{HOME}/devs/c++/clang-p2996/build/bin/clang",
"cpp26_lib":  "$ENV{HOME}/devs/c++/clang-p2996/build/lib",
\end{lstlisting}

\field{variables} declares a name for the shared part, which \code{\$NAME} (or
\code{\$\{NAME\}}) then stands for:

\begin{lstlisting}[style=json]
"variables": [
    { "name": "P2996", "value": "$ENV{HOME}/devs/c++/clang-p2996" }
],

"cpp26_root": "$P2996/build",
"cpp26_cxx":  "$P2996/build/bin/clang++",
"cpp26_cc":   "$P2996/build/bin/clang",
"cpp26_lib":  "$P2996/build/lib",
\end{lstlisting}

The substitution runs over \textbf{every string in the manifest} before any
field is read, so a variable works anywhere a string does ---
\field{user\_include}, \field{out\_dir}, a \field{user\_lib} directory, a
\field{compile\_definitions} value --- not only in the toolchain paths above.
The field itself disappears once applied.

Declarations are an \textbf{array}, not an object, because they are
\textbf{ordered}: a value may use the variables declared before it, and only
those.

\begin{lstlisting}[style=json]
"variables": [
    { "name": "TC",  "value": "/opt/toolchains/clang-p2996" },
    { "name": "BIN", "value": "$TC/build/bin" }
],
"cpp26_cxx": "$BIN/clang++"
\end{lstlisting}

\subsection{What is not substituted}

A manifest already carries two other dollar notations, and both belong to
\textbf{CMake}, which expands them at its own configure time. Rosetta hands
them over untouched, so the rule is that \textbf{only a name you declared is
substituted}:

\begin{center}
\small
\begin{tabular}{@{}lL{8.4cm}@{}}
\toprule
\textbf{Written} & \textbf{Result} \\
\midrule
\code{\$ENV\{HOME\}} & passed through verbatim --- CMake's environment lookup,
  whatever \field{variables} says \\
\code{\$\{CMAKE\_SOURCE\_DIR\}} & passed through verbatim --- no
  \code{CMAKE\_SOURCE\_DIR} was declared \\
\code{\$P2996} with \code{P2996} declared & substituted \\
\bottomrule
\end{tabular}
\end{center}

Declaring \code{ENV} is an error, since \code{\$ENV\{\ldots\}} could never
resolve to it.

\begin{warn}
The price of that leniency is that a \textbf{misspelt variable is silent}.
\code{\$P2996x} is one name, \code{P2996x}, which nothing declares --- so it
stays as written and surfaces later as a path that does not exist.
\code{\$\{P2996\}x} is how a variable abuts other text, and is worth preferring
for that reason.
\end{warn}

A variable name is a C identifier: letters, digits and \code{\_}, not starting
with a digit.

\section{Where the toolchain lives}
\label{sec:toolchain-paths}

Four optional fields point at the C++26 / P2996 reflection compiler used to
build the generator. Omit them all and rosetta uses these defaults:

\begin{center}
\small
\begin{tabular}{@{}ll@{}}
\toprule
\textbf{Field} & \textbf{Default} \\
\midrule
\field{cpp26\_root} & \code{\$ENV\{HOME\}/devs/c++/clang-p2996/build} \\
\field{cpp26\_cxx}  & \code{\$\{cpp26\_root\}/bin/clang++} \\
\field{cpp26\_cc}   & \code{\$\{cpp26\_root\}/bin/clang} \\
\field{cpp26\_lib}  & \code{\$\{cpp26\_root\}/lib} \\
\bottomrule
\end{tabular}
\end{center}

If your clang-p2996 build lives elsewhere, set \field{cpp26\_root} alone --- the
other three derive from it. Override an individual one only when the compiler
binaries or the \code{libc++} / \code{libc++abi} directory sit outside the usual
\code{bin/} and \code{lib/} layout.

\code{\$ENV\{HOME\}} is expanded by CMake at configure time, so the path stays
portable across machines. Of the targets, these fields affect only \lang{rest}
and \lang{json} --- every other one builds with an off-the-shelf compiler and
ignores them entirely. (The \emph{generator} always needs them.)

The command line wins over the manifest:
\code{rosetta\_gen --build --clang-p2996-root PATH}.

\field{qt\_dir} plays the same role for the \lang{qt} and \lang{qml} backends,
with \code{--qt-dir} as its flag.

\section{Naming things}

\begin{description}[leftmargin=0cm,style=nextline,font=\normalfont\code]
  \item[generator\_name] Basename of the emitted reflection driver:
    \code{"my\_person\_gen"} produces \code{gen/my\_person\_gen.cpp}. The built
    binary is always called \code{generator}, so this is mostly cosmetic --- its
    one load-bearing role is being the last-resort default for
    \field{module\_name}. Defaults to the manifest's directory name.
  \item[module\_name] The default binding module name, used by any bare-string
    target. A preset may supply one, which is why a preset-based manifest often
    names neither this nor \field{generator\_name}.
\end{description}

\begin{note}
The \emph{derived} \field{generator\_name} is folded to an identifier
(\code{my-cool-lib} becomes \code{my\_cool\_lib}) so a directory name cannot
produce an invalid module name. A value you spell out is used verbatim ---
if you write something that is not an identifier, you get it.
\end{note}

\section{Shared defaults: \field{namespace} and \field{header\_dir}}
\label{sec:shared-defaults}

When every entry repeats the same namespace and the same folder:

\begin{lstlisting}[style=json]
"classes": [
  {"name": "stressinv::Serie",      "header": "stressinv/Serie.h"},
  {"name": "stressinv::Data",       "header": "stressinv/Data.h"},
  {"name": "stressinv::CostMetric", "header": "stressinv/cost.h"}
]
\end{lstlisting}

factor them out:

\begin{lstlisting}[style=json]
"namespace": "stressinv",
"header_dir": "stressinv",
"classes": [
    {"header": "Serie.h"},
    {"header": "Data.h"},
    {"name": "CostMetric", "header": "cost.h"}
]
\end{lstlisting}

(\code{Serie} and \code{Data} also drop their \field{name}, which defaults from
the header stem --- and the derived name is namespace-qualified too.)

The rules apply identically to \field{classes}, \field{functions} and
\field{extensions}:

\begin{itemize}
  \item \field{namespace} qualifies an entry name \textbf{only when it carries
    no \code{::} of its own}. \code{"Serie"} becomes \code{stressinv::Serie};
    a name containing \code{::} passes verbatim. So fully qualified spellings,
    nested classes (\code{stressinv::Model::Inner}) and sub-namespaces keep
    working unchanged, and mixing shortened and full spellings in one manifest
    is fine.
  \item A \textbf{leading \code{::}} pins an entry to the global namespace:
    \code{"::c\_entry"} becomes \code{c\_entry}. This is the escape hatch for
    the odd global function (an \code{extern "C"} one, say) in an otherwise
    namespaced manifest.
  \item \field{header\_dir} is prepended to every entry header, inserting a
    \code{/} if missing: \code{"Serie.h"} becomes
    \code{stressinv/Serie.h}. A header living elsewhere can step out relative to
    it (\code{"../other/x.h"}), or you can leave \field{header\_dir} unset and
    spell every path in full.
\end{itemize}

\subsection{Grouped entries}
\label{sec:groups}

One pair of top-level defaults cannot cover a library whose headers live in
several folders, or that uses sub-namespaces. For that, an element of
\field{classes} or \field{functions} may be a \textbf{group}: an object
carrying \field{entries} plus its own local defaults, instead of being an entry
itself.

\begin{lstlisting}[style=json]
"namespace": "arch",
"header_dir": "Arch",
"classes": [
  {"name": "Vector3", "header": "math/math.h"},

  {"header_dir": "solvers", "entries": [
    {"name": "GmresSolver", "header": "Gmres.h"},
    {"header": "ParallelSolver.h"}
  ]},

  {"header_dir": "algos", "entries": [
    {"header": "DataSuperposition.h"},
    {"namespace": "sinv", "header_dir": "stress-inversion", "entries": [
      {"header": "JointData.h"},
      {"header": "types.h", "entries": [
        {"name": "MCMCConfig"},
        {"name": "MCMCResult"}
      ]}
    ]}
  ]}
]
\end{lstlisting}

\begin{itemize}
  \item A group's \field{header\_dir} \textbf{appends below} the inherited one:
    \code{Arch} + \code{solvers} gives \code{Arch/solvers/Gmres.h}.
  \item A group's \field{namespace} \textbf{appends to} the inherited one:
    \code{arch} + \code{sinv} gives \code{arch::sinv::JointData}. A leading
    \code{::} makes it absolute instead of appending.
  \item A group may set a \field{header}: the default header for entries that
    spell none. That is the natural shape for a run of classes declared by one
    header --- \code{types.h} above yields \code{arch::sinv::MCMCConfig} and
    \code{arch::sinv::MCMCResult}, both from
    \code{Arch/algos/stress-inversion/types.h}. Such entries need an explicit
    \field{name}, since the stem fallback would give every one the same name.
  \item Groups \textbf{nest}, \textbf{mix freely with plain entries} in the
    same array, and work identically under \field{functions} --- where a
    shared-header group reads especially well (a dozen shape generators all
    declared by \code{shapes.h}).
  \item A group cannot carry a \field{name}, and \code{//}-comment keys are
    ignored on groups as everywhere else.
\end{itemize}

Everything is resolved at load time, so backends and generated output are
byte-identical to the fully spelled form.

\section{Multiple include directories}

When your headers do not all live under one root, give \field{user\_include} an
array:

\begin{lstlisting}[style=json]
"user_include": ["./geom", "../shared/include", "/opt/thirdparty/include"]
\end{lstlisting}

Each entry follows the same resolution rules as the single-string form. Every
directory is added to the generated bindings' include path, so a class
\field{header} is resolved against \emph{all} of them --- the first directory
containing the file wins, exactly like a compiler's \code{-I} search order. The
array must not be empty; the single-string form is just the one-directory
shorthand.

\section{A field map}

The rest of Part~II is organised by what you are trying to do, rather than
alphabetically. Appendix~\ref{app:manifest} is the alphabetical version.

\begin{center}
\small
\begin{tabular}{@{}L{5.6cm}l@{}}
\toprule
\textbf{You want to\ldots} & \textbf{See} \\
\midrule
say which classes and functions to bind &
  Chapter~\ref{ch:whattobind} \\
choose the languages, and name the modules &
  Chapter~\ref{ch:targets} \\
get the generated projects to actually compile and link &
  Chapter~\ref{ch:building} \\
make a foreign container or an Eigen type cross the boundary &
  Chapter~\ref{ch:types} \\
reuse a fixed binding surface across projects &
  Chapter~\ref{ch:presets} \\
add docstrings, ranges, UI hints &
  Chapter~\ref{ch:annotations} \\
\bottomrule
\end{tabular}
\end{center}
